% =============================
% FAZIT & AUSBLICK
% =============================
\refstepcounter{chapter}
\chapter*{\thechapter. Fazit}
\addcontentsline{toc}{chapter}{\thechapter.\space Fazit \& Ausblick}

\section{Zusammenfassung}
Diese Arbeit hat gezeigt, dass Quantencomputer eine erhebliche Bedrohung für die Sicherheit aktueller Verschlüsselungsverfahren darstellen.
Solange die praktische Entwicklung von Quantencomputern unter der kritischen Grenze von etwa 4000 logischen Qubits verbleibt, bleibt der RSA grundsätzlich sicher. 
Allerdings zeigt die aktuelle Roadmap von IBM, dass diese Grenze in den 2030er Jahren erreicht werden könnte.
Die entscheidende Erkenntnis ist jedoch: Der Handlungsbedarf besteht bereits heute, nicht erst bei Marktreife. 
Besonders kritisch sind dabei langfristig vertrauliche Daten, die heute gesammelt und gespeichert werden, um sie in der Zukunft zu entschlüsseln.
\\
Die vom NIST standardisierten PQC-Verfahren bieten einen bewährten Rahmen für diese Transition. 
Besonders wichtig ist, dass diese Algorithmen bereits heute in kritischen Infrastrukturen wie Bankensystemen, Gesundheitswesen und Regierungskommunikation implementiert werden müssen. 
Harvest now, decrypt later--Angriffe unterstreichen die Dringlichkeit, da Daten, die heute mit RSA verschlüsselt werden, mit einem zukünftigen Quantencomputer mit 4000 logischen Qubits ausgelesen werden können.
Daher müssen PQC-Lösungen schon jetzt eingeführt werden. 

\section{Beantwortung der Fragestellung}
Die in dieser Arbeit untersuchte Fragestellung lautete: \textit{„Welche Bedrohung stellen Quantencomputer für die Sicherheit aktueller Verschlüsselungsverfahren dar, und inwiefern wären dadurch Daten gefährdet, wenn morgen ein Quantencomputer marktreif wäre?"}
\\
Ein Quantencomputer mit über 4000 logischen Qubits würde die mathematischen Grundlagen der heute verwendeten asymmetrischen Verschlüsselungsverfahren vollständig aufheben.
Dadurch wäre es möglich, alle damit verschlüsselten Daten innerhalb kurzer Zeit zu entschlüsseln. %Zudem besteht die Gefahr eines sogenannten Harvest now, decrypt later-Angriffs. 
%Dies betrifft unmittelbar persönliche Daten wie Finanztransaktionen und Bankinformationen, medizinische Daten, private Kommunikation und Authentifizierungsdaten. 
\\
Bei Quantencomputer mit 2000 logischen Qubits, welcher für Forschungs- und Industriezwecke ausreichen sollte, wäre die gängigsten Verschlüsselungsverfahren noch sicher. 
Allerdings könnten gezielte Angriffe auf schwächere Systeme oder Protokolle möglich sein, was die Notwendigkeit unterstreicht, frühzeitig auf quantensichere Algorithmen umzustellen.

\section{Kritische Reflexion}
Wie jede technologische Entwicklung haben auch Quantencomputer Chancen und Risiken.
Während sie die Sicherheit bestehender Verschlüsselungsverfahren bedrohen, bieten sie gleichzeitig die Möglichkeit, 
für mathematische Beweise und effiziente Algorithmen für komplexe Probleme.
Der Zeitpunkt des Q-Day ist ungewiss. Damit bleiben viele Fragen offen. 
Auch können technologische Durchbrüche die Entwicklung beschleunigen, staatliche Regulierungen beispielsweise könnten sie hingegen verzögern.
\\
Selbst Fachpersonen verstehen noch nicht alle Verhaltensmerkmale der quantenphysikalischen Welt. Viel Wissen und Fachpersonen mit diesem Wissen müssen aufgebaut werden.
Der vorgegebene Umfang dieser Berufsmaturitätsarbeit erlaubte nur einen kleinen Einblick in die komplexe Thematik der Quantencomputer und deren Auswirkungen auf die Datensicherheit.

